% !TEX root = ../notes.tex

\documentclass[../notes.tex]{subfiles}

\begin{document}

\section{April 15}
Here we go.

\subsection{Preorientations}
We are still interested in deforming some $p$-divisible groups. For example, given a connective $\mathbb E_\infty$-ring $R$ for which $R\coloneqq\tau_{\ge0}A$ where $A$ is weakly periodic, we have the $p$-divisible group
\[\op{Spf}\tau_{\ge0}\op{Hom}_\mathbb S\left(\Sigma_+^\infty\mathrm B\ZZ/p^n\ZZ,A\right).\]
When $R$ is classical, our homotopy-coherent definitions recover the classical notions, but sometimes we get more interesting things.
\begin{example}
	Suppose $E$ is an elliptic curve over a classical ring $R$ of positive characteristic $p$. Then there is the $p$-divisible group
	\[E[p]\into E\left[p^2\right]\into E\left[p^3\right]\into\cdots.\]
	This $p$-divisible group is connected if and only if $E$ is supersingular. Thus, it typically does not come from a formal group.
\end{example}
\begin{example}
	Over any connective $\mathbb E_\infty$-ring $R$, we have $\widehat{\mathbb G}_m=\op{Spf}\left(R\otimes\Sigma^\infty_+\ZZ\right)^\land_{t-1}$. (Intuitively, we should think about this as $R[t,1/t]$ completed at the identity $t=1$.) It turns out that $\mathbb S$ is the universal deformation of $\widehat{\mathbb G}_m$ over $\ZZ$: namely, $\mathbb S$ is the groupoid of formal deformations of $\widehat{\mathbb G}_m$, which we defined earlier. We may say that $(\mathbb S,\widehat{\mathbb G}_m)$ is the universal deformation.
\end{example}
Even though the above toy example might not seem so interesting, it will allow us to recover $\mathrm{KU}$.
\begin{definition}[preorientation]
	Fix an $\mathbb E_\infty$-ring $R$. Then a \textit{preorientation} of a formal group $\widehat{\mathbb G}$ over $\tau_{\ge0}R$ is a map $\mathbb S^2\to\widehat{\mathbb G}(\tau_{\ge0}R)$ of pointed spaces.
\end{definition}
\begin{remark}
	If $R$ is weakly periodic, then the space of preorientations on $\widehat{\mathbb G}$ is
	\[\op{Hom}_{\mathrm{Spaces}_*}\left(\mathbb S^2,\Omega^\infty\widehat{\mathbb G}(\tau_{\ge0}R)\right)=\op{Hom}_{\mathrm{Spaces}_*}\left(\Sigma^2\ZZ,\widehat{\mathbb G}(\tau_{\ge0}R)\right).\]
	We claim that this is the same as maps from the Quillen formal group $Q\coloneqq\op{Spf}\tau_{\ge0}\op{Hom}\left(\Sigma^\infty_+\CP^\infty,R\right)$. Indeed, this follows by some unwinding once we note $\Sigma_+^\infty\CP^\infty=\Sigma_+^\infty\Omega^\infty\Sigma^2\ZZ$.
\end{remark}
The perspective from the Quillen formal group motivates the following definition.
\begin{definition}[orientation]
	Fix a weakly periodic $\mathbb E_\infty$-ring $R$. An \textit{orientation} is a preorientation such that the induced map from the Quillen formal group is an equivalence.
\end{definition}
\begin{example}
	The space of preorientations of $\widehat{\mathbb G}_m$ over $\mathbb S$ is
	\[\op{Hom}_{\mathrm{Spaces}_*}\left(\mathbb S^2,\widehat{\mathbb G}_m(\mathbb S)\right).\]
	The target is $\op{Hom}_{\mathrm{Ring}}(\Sigma^\infty_+\Omega^\infty\ZZ,\mathbb S)$, and there is an adjunction to $\op{Hom}_\mathbb S(\ZZ,\mf{gl}_1(\mathbb S))$, so this is the same as $\op{Hom}_\mathbb S\left(\Sigma^2\ZZ,\mf{gl}_1(\mathbb S)\right)$. This is the same as $\mathbb E_\infty$-maps $\Sigma_+^\infty\Omega^\infty\Sigma^2\ZZ\to\mathbb S$, which is the same as maps $\Sigma_+^\infty\CP^\infty\to\mathbb S$. This explains the terminology. While we're here, we note that 
\end{example}
\begin{remark}
	The moral of the previous example is that we can recover $\mathrm{KU}$ by inverting the Bott element $\beta$ in $\Sigma^\infty_+\CP^\infty$, which we can kind of view as the space of preorientations of $\widehat{\mathbb G}_m$.
\end{remark}
Anyway, we see that a preorientation is some sort of general notion of an orientation which works even when the $\mathbb E_\infty$-ring is not $\tau_{\ge0}R$ of an even $\mathbb E_\infty$-ring $R$.

\subsection{The Bott Map}
Now, we see we are interested in determining when the induced map from the Quillen formal group is an equivalence. Intuitively, we are kind of asking for a map $k[[x]]\to k[[x]]$ induced by a formal power series to be an isomorphism, and we remark that we can basically check if this is an isomorphism on $I/I^2$, where $I=(x)$. (This is some Nakayama's lemma argument.) To use this idea, it will be useful to make the following definition.
\begin{definition}
	Fix an $\mathbb E_\infty$-ring $R$ and a $1$-dimensional formal group $\widehat{\mathbb G}$ over $\tau_{\ge0}R$. Then we define
	\[\omega_{\widehat{\mathbb G}}\coloneqq R\otimes_{\OO_{\widehat{\mathbb G}}}\op{fib}\varepsilon,\]
	where $\varepsilon$ is the identity $\widehat{\OO}_{\widehat{\mathbb G}}\to R$.
\end{definition}
\begin{remark}
	It turns out that $\omega$ is related to some kind of cotangent complex.
\end{remark}
\begin{remark}
	For dimension reasons, $\omega$ is an invertible $R$-module of rank $1$.
\end{remark}
\begin{remark}
	It turns out that $\omega$ fits in a fiber sequence $\Sigma\omega\to R\otimes_{\widehat{\OO}_{\widehat{\mathbb G}}}R\stackrel m\to R$. This is analogous to the fiber sequence
	\[\Sigma k\to k\otimes_{k[[x]]}k\to k.\]
	% More formally, the moral is that $\widehat{\OO}_{\mathbb G}$ fits into the following diagram
	% % https://q.uiver.app/#q=WzAsNCxbMCwwLCJcXHdpZGVoYXR7XFxPT31fe1xcd2lkZWhhdHtcXG1hdGhiYiBHfX0iXSxbMSwwLCJSIl0sWzAsMSwiUiJdLFsxLDEsIlJcXG90aW1lc197XFx3aWRlaGF0e1xcT099X3tcXHdpZGVoYXR7XFxtYXRoYmIgR319fVIiXSxbMCwxXSxbMSwzXSxbMiwzXSxbMCwyXV0=&macro_url=https%3A%2F%2Fraw.githubusercontent.com%2FdFoiler%2Fnotes%2Fmaster%2Fnir.tex
	% \[\begin{tikzcd}[cramped]
	% 	{\widehat{\OO}_{\widehat{\mathbb G}}} & R \\
	% 	R & {R\otimes_{\widehat{\OO}_{\widehat{\mathbb G}}}R}
	% 	\arrow[from=1-1, to=1-2]
	% 	\arrow[from=1-1, to=2-1]
	% 	\arrow[from=1-2, to=2-2]
	% 	\arrow[from=2-1, to=2-2]
	% \end{tikzcd}\]
	% where we view the bottom-left $R$ as the free $\mathbb E_\infty$-algebra over the point.
\end{remark}
Let's do a calculation to see $\omega$ appear.
\begin{lemma}
	Fix a connective $\mathbb E_\infty$-ring $R$ and a $1$-dimensinal formal group $\widehat{\mathbb G}$. For any $\mathbb E_\infty$-algebra $A$, there is a natural map
	\[\Omega^2\widehat{\mathbb G}(A)\to\op{Hom}_R\left(\omega,\Sigma^{-2}A\right).\]
\end{lemma}
\begin{proof}
	We compute directly. Note that $\Omega\widehat{\mathbb G}(A)$ is the same as a map $R\otimes_{\OO_{\widehat{\mathbb G}}}R\to A$ of $\mathbb E_\infty$-algbras. This embeds into maps of $R$-modules, and then it restricts to $R$-module maps $\Sigma\omega\to A$. This space is the same as $\Omega\mathrm{Hom}_R(\omega,A)$, as desired.
\end{proof}
\begin{example}
	Now, if $A$ is weakly periodic, then let $Q$ be the associated formal group. Then $\Sigma\omega_Q\to A\otimes_{\op{Hom}(\Sigma_+^\infty\CP^\infty,A)} A\to A$ has its middle term given by
	\[\op{Hom}_{\mathbb S}(\Sigma_+^\infty(\mathrm{pt}\times_{\CP^\infty}\mathrm{pt},A).\]
	Now, the left space is basically $S^1$, which we see by computing this fiber product. Thus, the total space of maps is $A\oplus\Sigma^{-1}A$, and it turns out that the induced map $A\oplus\Sigma^{-1}A\to A$ is the left projection. Thus, $\omega_Q=\Sigma^{-2}A$.
\end{example}
The moral of this calculation is that we do have an isomorphism $\omega_Q\to\Sigma^{-2}R$ whenever $R=\tau_{\ge0}A$ for a weakly periodic $A$, and a preorientation gives the choice of any such map $\omega_Q\to\Sigma^{-2}R$, and it upgrades to an isomorphism if and only if the preorientation was actually an isomorphism all along.

The following construction came out of these computations.
\begin{definition}[Bott map]
	Fix a formal group $\widehat G$ over a connected $\mathbb E_\infty$-ring $R$. Given a preorientation, the \textit{Bott map} is the induced map $\omega_{\widehat G}\to\Sigma^{-2}R$.
\end{definition}
\begin{remark}
	In general, the Bott map is an $R$-module map of invertible $R$-modules of rank $1$, but it does not have to be an isomorphism.
\end{remark}
\begin{example}
	Applied to $\widehat{\mathbb G}_m$ over $\mathbb S$, we find that $\omega=\Sigma_+^\infty\CP^\infty$. In this case, the Bott map above turns out to be induced by $\beta$.
\end{example}

\subsection{Morava \texorpdfstring{$E$}{ E}-Theory}
Here is our definition.
\begin{definition}[Morava]
	Fix a pair $(R,\Gamma)$, where $R$ is a classical $\FF_p$-algebra, and $\Gamma$ is a $p$-divisible group. Suppose further that $L_R^{\mathrm{alg}}$ is almost perfect, and $\Gamma$ is non-stationary. Then we define $R^{\mathrm{def}}$ via \Cref{thm:lurie} to be the universal $\mathbb E_\infty$-algebra with a preorientation (of its underlying formal group). We then define $E(R,\Gamma)$ to be the $\mathbb E_\infty$-ring to be the localization of $R^{\mathrm{def}}$ at the Bott element, meaning that it is the colimit of
	\[R\to\omega^{-1}\to\omega^{-2}\to\cdots.\]
\end{definition}
\begin{remark}
	We have not shown that $R^{\mathrm{def}}$ actually exists. The point is that we add to the deformation problem $R_0^{\mathrm{def}}$ the data of the map from $\mathbb S^2$. Note that $R^{\mathrm{def}}$ is an algebra over $R_0^{\mathrm{def}}$.
\end{remark}
\begin{remark}
	This is called $E$-theory after Morava's wife Eleanor.
\end{remark}
\begin{definition}[Luriean]
	Fix a pair $(R,\Gamma)$, where $R$ is a classical $\FF_p$-algebra, and $\Gamma$ is a $p$-divisible group. Then $(R,\Gamma)$ is \textit{Luriean} if and only if the following hold.
	\begin{listalph}
		\item The module $L_R^{\mathrm{alg}}$ is almost perfect.
		\item The group $\Gamma$ is non-stationary.
		\item The map $R_0^{\mathrm{def}}\to E(R,\Gamma)$ is an isomorphism on $\pi_0$.
		\item The odd homology of $E(R,\Gamma)$ vanishes.
	\end{listalph}
\end{definition}
Here are some examples about some pairs $(R,\Gamma)$ which are Luriean.
\begin{example}[topological modular forms]
	If $A$ is an elliptic curve over an $\mathbb F_p$-algebra $R$, then we can build $E(R,A\left[p^\infty\right])$. Taking the limit, we receive the $\mathbb E_\infty$-ring
	\[\mathrm{TMF}^\land_p=\lim_{\substack{A\colon\Spec R\to\mc M_{\mathrm{ell}}\\A\text{ \'etale}}}E\left(R,A\left[p^\infty\right]\right).\]
	This is interesting to mathematical physicists.
\end{example}
\begin{theorem}[Goerss--Hopkins--Miller]
	Fix a perfect field $K$ of positive characteristic $p$, and choose a one-dimensional formal group $\widehat{ G}_0$ of finite height $n$. Then $E(K,\widehat G_0)$ exists, and
	\[\pi_*E(K,\widehat G_0)=W(k)[[v_1,\ldots,v_{n-1}]][u,1/u].\]
	Here, $W(k)$ are the Witt vectors, $\deg v_i=2p^i-2$ (coming from $\mathrm{BP}$!), and $\deg u=2$.
\end{theorem}
Here, recall that the height of a $1$-dimensinal formal group is the $n$ for which $[p](x)=x^{p^n}$ for some choice of coordinate $x$. For example, the height classifies formal groups over an algebraically closed field.
\begin{remark}
	Set $E\coloneqq E(K,\widehat G_0)$, and suppose that $E$ is $K(n)$-local, and $A$ is $K(n)$-local and even periodic. Then
	\[\op{Hom}_{\mathbb E_\infty}(E,A)=\op{Hom}((K,\widehat G_0),(\pi_0A/(p,v_1,\ldots,v_{n-1}),G_Q)),\]
	where $G_Q$ is the Quillen formal group. Amazingly, this target is discrete (because the target is morphisms of classical formal groups).
\end{remark}
Here is another result.
\begin{definition}[Nullstellensatzian]
	A $K(n)$-local $\mathbb E_\infty$-ring $A$ is \textit{Nullstellensatzian} if and only if every compact object in the category of $K(n)$-local $\mathbb E_\infty$-algebras maps to $A$.
\end{definition}
\begin{remark}
	In classical algebraic geometry over a field $k$, we see that the compact objects are the algebras of finite type, and the Nullstellensatz says that all such objects admit maps to the algebraic closure $\overline k$.
\end{remark}
\begin{theorem}
	The collection of nonzero Nullstellensatzian rings are exactly the Morava $E$-theories $E(K,\widehat G_0)$, where $K$ is algebraically close, and $\widehat G_0$ has a finite height.
\end{theorem}
Thus, the Morava $E$-theories behave like algebraically closed fields.

\end{document}